\nwuFrontHeading{ABSTRACT}

Platformized communication has reshaped the career transition of public entertainers. Print magazines, screen productions, variety programs, commercial endorsements, and self-operated social media channels are no longer merely separate exposure outlets; together, they form a dynamic system of professional identity, narrative rhythm, and public relationship management. This public demonstration uses the career trajectory of Nozomi Sasaki as a case to examine how a Japanese female entertainer moves across print media, screen industries, brand communication, and creator-led video platforms while maintaining visibility and image resilience after public controversy.

The thesis first establishes a factual baseline from official profiles and public news sources: Sasaki was born on February 8, 1988, in Akita Prefecture, entered the entertainment industry in 2006, has continued to work across films, television dramas, commercials, and magazines, and opened the official YouTube channel ``Nozomi Sasaki (tentative)'' in 2024. The thesis then organizes the career path into four analytical stages: print-based recognition, screen and commercial expansion, family-related public controversy, and platform-led restart. The analysis focuses on platform control, content rhythm, role boundary, and crisis response. Reused public demo images and tables are included to verify how a doctoral thesis template handles figures, cross references, equations, references, appendix materials, achievements, acknowledgements, and author biography.

This sample is not an entertainment-news digest and does not make moral judgments about private life. All factual claims are based on public webpages and news reports, and sensitive events are mentioned only as background for analyzing how public narratives affect career trajectories. At the typesetting level, the project preserves the required structure of \texttt{thesis-nwu-doctor}, including the Chinese cover, English title page, intellectual property statement, originality statement, bilingual abstracts, table of contents, chapters, conclusion, references, appendix, doctoral achievements, acknowledgements, and author biography.

\nwuKeywordsEnLine
